\documentclass[a4paper,12pt]{article}

\usepackage{graphicx}
\usepackage{amsmath}
\usepackage{float}

\begin{document}

\noindent
\begin{minipage}[t]{0.2\textwidth}

\includegraphics[width=2.5cm]{iiitb_logo.png}
\end{minipage}
\hfill
\begin{minipage}[t]{0.75\textwidth}
\raggedleft

{\Large \textbf{Pavithra Kune}}

COMETFWC069

\end{minipage}

\vspace{0.8cm}

\begin{center}
{\Large \textbf{GATE QUESTION Q36}}
\end{center}

\vspace{0.3cm}

\hrule

\vspace{0.5cm}

\begin{figure}[H]
\centering

\includegraphics[width=0.8\textwidth]{question.png}
\caption{GATE Question Q36}
\end{figure}

\tableofcontents

\newpage

\section{Components}

\subsection{Breadboard}

A breadboard is used to make temporary circuit connections.

\subsection{Seven Segment Display}

A common anode seven segment display is used to display the output.

\subsection{Arduino Uno}

Arduino Uno is used to implement the Boolean expression.

\subsection{Jumper Wires}

Jumper wires are used to connect the circuit components.

\section{Boolean Simplification}

Given,

[
[X+Z{\bar Y+(\bar Z+X\bar Y)}]
[\bar X+\bar Z(X+Y)]
]

For (X=1),

[
F=\bar Z
]

For output 1,

[
Z=0
]

\section{Hardware Connections}

\begin{center}
\begin{tabular}{|c|c|}
\hline
Arduino Pin & Segment\
\hline
D2 & a\
D3 & b\
D4 & c\
D5 & d\
D6 & e\
D7 & f\
D8 & g\
\hline
\end{tabular}
\end{center}

The common anode pin is connected to +5V.

\section{Software Implementation}

\begin{verbatim}
int Z = 0;
int F = !Z;

if(F == 1)
{
display1();
}
else
{
display0();
}
\end{verbatim}

\section{Output}

\begin{figure}[H]
\centering
\includegraphics[width=0.5\textwidth]{output1.jpg}
\caption{Output showing 1}
\end{figure}

\begin{figure}[H]
\centering
\includegraphics[width=0.5\textwidth]{output0.jpg}
\caption{Output showing 0}
\end{figure}

\section{Result}

The given Boolean expression was implemented using Arduino Uno and a seven segment display.

[
F=\bar Z
]

Hence,

[
Z=0
]

for output 1.

\end{document}
